\section{Centered heights and Haar densities}\label{sec:heights}

The local arguments have supplied upper bounds for the missing panel gaps
and positive heights at isolated types. We first use those isolated
heights to prove tightness of their centered gaps. Together with the panel
upper bounds, this permits a comparison of source Haar volume and target
lattice packing on paths selecting one vertex in each component. The
comparison fixes the panel speeds and makes all these tuple gaps tight.
We then construct the remaining heights from their centered laws.

The resulting pair identities identify the full centered Cartan law.
A second, angularly localized packing argument shows that inverse chamber
heights turn the boundary measurements into ordinary chamber probability.
The associated isometry and its exact height intertwining are the inputs
to the common graph domains in Section~\ref{sec:locality} and to the
real-pencil argument in Section~\ref{sec:gluing}.

\subsection{Paths, line norms, and local-field conventions}

For every non-isolated source position $i$, put
\[
 s_i=\frac{d_S}{d_T},
\]
where $S,T$ are its source and matched target panel groups and the constants
are those of Sections~\ref{sec:finite} and~\ref{sec:panels}.  At an isolated
position put $s_i=1$.  At non-isolated positions the constants satisfy the
upper-speed estimates \eqref{eq:finite-speed} and \eqref{eq:input-11}
at finite and archimedean places, respectively.  All these constants have
opposition symmetry $s_i=s_{i^*}$ and product one on every cycle of the
type permutation $p$: the ratios of the $d_R$ telescope, and the
isolated constants are one.

We use mixed lexicographic paths with one missing high type per indicated
component and all isolated types retained.  A further completing translation
is applied to the inputs before the faster limit is taken.  The sharp face
marginals are unchanged by translations stabilizing those faces, although
the residual projectivity may change.  Adding a missing direction strictly
refines the relevant inner stratum, by the gap-count argument of
Section~\ref{sec:links}.  Completing the chamber makes every target gap
diverge.  Repeating a selected direction by a slower translation preserves
its image type and cannot add a type beyond the upper count.

Compact face frames may be chosen by Borel compact-transport sections.
When a frame function is passed through a link or through a completing
limit, first restrict to its compact continuity sets for the fixed input
laws.  Those laws are unchanged by the translations stabilizing the sharp
faces.  Their complements have arbitrarily small mass on each fixed
testing vector.  Countably many determining vectors and charts, followed
by strong continuity, therefore justify the framed substitutions without
assuming continuity of a global Borel section.

Offsets will always be held fixed during every ordinary path time and
exhausted only afterwards.  Thus even a negative offset in a selected
positive direction does not undo that direction's divergence.  The parameter
$t_i$ denotes its total value, including such offsets.  At a finite place
all split translations and offsets run on the fixed sufficiently divisible
coweight mesh of Section~\ref{sec:setup}.

\begin{table}[!htbp]
\centering
\footnotesize
\setlength{\tabcolsep}{3pt}
\renewcommand{\arraystretch}{1.15}
\begin{tabular}{@{}>{\raggedright\arraybackslash}p{22mm}>{\raggedright\arraybackslash}p{58mm}>{\raggedright\arraybackslash}p{57mm}@{}}
\toprule
Local type & Angular measurements and gap estimates & Covariance and domains \\
\midrule
Finite-place panels
\newline Section~\ref{sec:finite}
&
Each endpoint pencil is a commuting family of evaluations of homogeneous
analytic homeomorphisms. The reduced output projectivity is $ab^{-1}$
(Proposition~\ref{prop:finite-maps}). On the permissible cocharacter mesh,
$d_T\ell_T(g)\le d_St+O_{\mathrm{tight}}(1)$, uniformly in slower
translations with the same total gap $t$
(Proposition~\ref{prop:finite-speed}).
&
The source acts through $\eta:S\to T$ with isomorphic differential
over the common prime field $\mathbb Q_p$
(Proposition~\ref{prop:finite-torsor}). The original endpoint observations
and exact input translations are retained. This branch uses bounded
spectral calculi. Its simultaneous upper-tail estimate uses a finite union
bound over coordinates; it requires no commutation between different
pencils.\\[5pt]
Non-isolated archimedean panels
\newline Section~\ref{sec:panels}
&
Sharp endpoint faces and missing-vertex measurements retain the ordered
residue graph relation throughout the successive completing limits.
The estimate is $d_T\ell(g)\le d_St+O_{\mathrm{tight}}(1)$ on every fixed
input, including the stated subsequent offset limits
(Proposition~\ref{prop:panels-upper}).
&
Source and measured-group covariance retain their compact-frame
corrections. The scalar section domain $\mathscr D$ is dense and
transport-invariant (Lemma~\ref{lem:panels-global-domain}). The
finite-dimensional algebra acts on its section space with the positive
form identity of Proposition~\ref{prop:panels-positive-action}; action
exchange is estimated on fixed domain vectors
(Lemma~\ref{lem:panels-action-exchange}).\\[5pt]
Isolated real factors
\newline Sections~\ref{sec:ray}--\ref{sec:scale}
&
The pure-vertex measurements are sharp, and complementary lengths are
tight. A subsequence carries uniformly positive mass in a linear window
on every fixed vector (Proposition~\ref{prop:scale-linear-subsequence}).
The integrated pair estimate gives $Q_S=Q_T=Q$ and square-integrable
horizontal derivatives (Lemma~\ref{lem:scale-equal-dimension}).
&
The closed derivative construction gives positive injective heights
$b_x$ strongly commuting with their own $P_x$. They satisfy
$\sigma_a(b_x)=\operatorname{dil}(a,x)^{-Q}b_{ax}$ and
$\beta_h(b_x)=\operatorname{dil}(h^{-1},P_x)^{-Q}b_x$.
Closure of the module core and normal copying are proved before use
(Lemma~\ref{lem:scale-module-extension} and
Proposition~\ref{prop:scale-isolated-height}).\\
\bottomrule
\end{tabular}
\caption{The three local inputs to product calibration. Here $P_z,Q_z$
are the endpoint projection-valued measures, $t$ is the total completing
source gap, $\ell$ is the residual target gap, $d_R$ is the comparison
constant for a non-isolated panel type, and $Q$ is the homogeneous
dimension of an isolated boundary. All tightness and linear-window
assertions concern fixed Hilbert vectors.}
\label{tab:local-branches}
\end{table}

At this stage, incident labels commute by
Proposition~\ref{prop:transfer-sharp-flags}; the stronger commutation of
arbitrary labels is still to be proved. The isolated linear-window
estimate becomes tightness of centered gaps in
Lemma~\ref{lem:heights-tightness-test}. The present section then supplies
the Haar-height isometry and exact height intertwining. The graph cores
needed for unbounded locality, and commutation between the two normal
copies at arbitrary isolated labels, are established in
Section~\ref{sec:locality}. Real-pencil commutation and simultaneous
chamber commutation follow in Section~\ref{sec:gluing}.


An event of positive infinite excess or deficit is unchanged by fixed left
or right target shifts: a fixed multiplication changes the Cartan vector
by a bounded amount.  The independent contracted invariances of the mixed
link therefore make its compressed mass scalar.  At a sharp angular slot
its angular subeffect is that scalar times the sharp marginal: its
compression is dominated by the relevant projection at every cut.  This
observation applies after retaining additional length or phase variables.
It permits us to fix a purported positive mass in an iterated limit, work
with bounded shift windows and continuous subtests with slack, and only
then exhaust thresholds.  No common good set of lattice indices for all
rotations and offsets is being asserted.

Let $\lambda_i$ be the weight of the nilradical determinant line at the
standard vertex of type $i$.  It is a positive multiple of the corresponding
fundamental weight, and
\[
 \lambda_i(e_i)=2\rho(e_i).
\]
Indeed the positive roots omitted from that nilradical involve only the
other simple roots and hence vanish on $e_i$.  Superscripts $S,T$ distinguish
the source and target lists.  Write $j(x,g)$ for the norm growth of the
determinant line at $x$ under $g$, and $k(y,x)$ for its normalized absolute
pairing with the dual line at the opposite-type vertex $y$.  The normalization
is one at a standard compact-frame opposite pair.  Consequently
\[
 k(gy,gx)j(y,g)j(x,g)=k(y,x).
\]
The product of line cocycles with total weight $2\rho$ will be denoted by
$j_{2\rho}$.  It is the inverse Jacobian of the compact chamber probability.

\begin{lemma}[Exact line normalization]\label{lem:heights-line-normalization}
Give the determinant line and its dual homogeneous norms obtained from
special-compact-invariant norms.  At a finite place use the module absolute
value, including its residue-degree normalization.  For an actual Cartan
centralizer representative $z$, the highest-line growth is exactly
\[
 |\chi(z)|=\exp\bigl(\lambda_i(\mu(z))\bigr).
\]
When the $i$-gap diverges, the normalized operator has only its highest-line
block in the limit.  If $w$ is its limiting reverse dual vertex and $v$ is
transverse to $w$, then
\begin{equation}\label{eq:input-14}
 j_T(v,h^{-1})\exp\bigl(-\lambda_{p(i)}^T(\mu(h^{-1}))\bigr)
 \longrightarrow k_T(w,v).
\end{equation}
The convergence is uniform with all needed derivatives on compact
transverse real patches.  In finite-place variables the normalized ratios
are eventually locally constant on compact transverse boxes with a uniform
level.  Fixed multiplications on the two sides change the logarithmic top
growth by the corresponding endpoint line cocycles.
\end{lemma}

\begin{proof}
Use the actual dominant translation vector of the centralizer, as fixed in
Section~\ref{sec:setup}; it records the exact module of each centralizer
character.  The remaining centralizer parts and their inverses are uniformly
bounded on the weight spaces.  Highest-weight strings separate every other
block from the highest line by a positive multiple of the pertinent gap.
After division by the exact highest-line growth those blocks therefore
vanish.  The compact factors are isometries for the chosen norms.

In compact frames the limit is the rank-one operator
$\widetilde u\otimes\widetilde w$, with unit lifts of the extremal lines,
up to a scalar of absolute value one.  Evaluating it on the line at $v$
gives its normalized pairing with $w$, proving
\eqref{eq:input-14}.  Multiplying on either side transforms the relevant
extremal line; the norm needed to renormalize that line is exactly its
$j$-factor.  At a finite place one can alternatively evaluate on a compact
transverse test line and divide by its reverse pairing.  Once the gap is
large, the ultrametric inequality makes the resulting norm identity exact.
All denominators have a uniform lower bound on the buffered patches, so the
locally constant tests have uniform levels there.  In particular there is
no additional bounded but nonvanishing Cartan error in
\eqref{eq:input-14}.

For the Jacobian assertion, reduce by compact frames to a minimal parabolic
element.  Its action on the opposite-root chart has determinant given by
the product of positive-root modules, including double roots and their
multiplicities.  This is $\exp(2\rho)$ and gives the inverse compact-angular
Jacobian.  At an isolated real vertex the horizontal dilation gives the same
Jacobian by its tangent grading.  Thus its $j$ is exactly the cocycle of the
height already constructed in Section~\ref{sec:scale}; equality of the
homogeneous dimensions in Proposition~\ref{prop:scale-isolated-height} yields
\[
 \lambda_i^S(e_i)=\lambda_{p(i)}^T(e_{p(i)})
 \qquad\text{at isolated positions}.
\]
\end{proof}

At finite places a continuous real logarithmic character of a parabolic
vanishes on its radical by contraction, on compact subgroups by compactness,
and on root-generated semisimple directions.  It therefore depends only on
the split center, up to compact or finite-index changes.  The selected
coweight mesh is a full lattice in those split-center coordinates and
determines such a character.  This is the character convention used below;
no real derivative on a building is needed.

\subsection{Moving tests and ordered transport}

We first record a uniform version of the smooth-test extension needed in
the transfer argument.  The coefficient traces and matrix sizes have the
normalization fixed in Proposition~\ref{prop:transfer-transverse}.  When
writing a constant from the transverse-product estimate, we retain the
homogeneous Hilbert bounds in \eqref{eq:transfer-homogeneous}, or its
trace-mass factor after an operator-norm-only bound.  In every application
below the cut arrays are made from a fixed bounded-column family, so those
factors are bounded.

\begin{lemma}[Radial tests with moving centers]\label{lem:heights-moving-tests}
On the buffered transverse supports of \eqref{eq:input-2}, one may add
finitely many tests
\[
 \varphi\!\left(\frac{\ell-b}{R}\right),
 \qquad \varphi\in C_c^\infty(\mathbb R^q),\quad
 b\in\mathbb R^q,\quad R\geq1,
\]
where $\ell$ is a fixed finite list of logarithmic matrix norms, including
inverse norms.  At real places use Hilbert--Schmidt matrix norms, which
are smooth away from zero; at finite places use the maximum norms fixed
in Section~\ref{sec:setup}.  One may also add characters of these logarithmic norms
with frequencies in a fixed bounded set.  The transverse-product bounds
are uniform in $b,R$ and those frequencies, with constants depending on
the fixed test functions and the original buffered supports.  The same
assertion holds in mixed real and finite-place coordinates, with compact
clopen tests in the latter.
\end{lemma}

\begin{proof}
For a length-increasing substitution $p\mapsto up$, normalized matrix
multiplication gives, componentwise,
\[
 \ell(up)=\ell(u)+\ell(p)+q(S_u,S_p).
\]
Here $S$ denotes projectively normalized matrix coordinates, and $q$ has a
smooth extension with uniform derivative bounds on the compact coordinate
boxes.  Its denominator is bounded below by transversality.  For inverse
norms use $(up)^{-1}=p^{-1}u^{-1}$, so the normalized product has the
reversed order.  At finite places $q$ is bounded and locally constant with
a uniform level on the same kind of buffered normalized boxes, by
Lemma~\ref{lem:heights-line-normalization} and the local-field conventions.

Fourier inversion gives
\begin{align*}
 \varphi\!\left(\frac{\ell(up)-b}{R}\right)
 =\int_{\mathbb R^q}\widehat\varphi(\xi)
 &\exp\!\left(\frac{\mathrm i\xi\cdot(\ell(u)-b)}{R}\right)
 \exp\!\left(\frac{\mathrm i\xi\cdot\ell(p)}{R}\right)\\
 &\hspace{8mm}\cdot
 \exp\!\left(\frac{\mathrm i\xi\cdot q(S_u,S_p)}{R}\right)\,d\xi,
\end{align*}
with the Fourier normalization absorbed in $\widehat\varphi$.
The first factor is a scalar of modulus one; it contains the unrestricted
center $b$.  The second is a bounded unary multiplier on the shorter
argument.  Its logarithmic argument is not restricted to a compact set and
is not differentiated as a rescaled compact coordinate.  Only the third
factor is expanded in the compact normalized coordinates.

For each fixed derivative order $r$, that compact factor has seminorm at
most $C_r(1+|\xi|)^r$, uniformly in $R\geq1$, the center, and the supported
shifts.  Smooth separated expansions on a fixed larger box therefore have
a summable supremum bound of that polynomial size.  In a finite-place slot
one uses the fixed finite partition into clopen boxes instead.  Apply the
unary Gram--Plancherel estimate \eqref{eq:transfer-homogeneous} at each
frequency, retaining any designated shorter cutoff outside the expansion.
Minkowski's inequality and
\[
 \int_{\mathbb R^q}|\widehat\varphi(\xi)|(1+|\xi|)^r\,d\xi<\infty
\]
give the required uniform bound.  The inequality holds because
$\widehat\varphi$ is Schwartz.  A bounded additional character frequency
replaces $\xi/R$ by $\theta+\xi/R$ in the compact factor, with the same
polynomial estimate.  Finitely many cutoffs use joint frequency variables
and products of Schwartz transforms.

The compact expansion uses a common unary basis with coefficient bound
$\sum_n\sup_u|c_n(u)|<\infty$, as in
\eqref{eq:transfer-unary-expansion}; a bound with the supremum and sum in
the opposite order would not suffice.  The frequency integral above
retains exactly this stronger bound.

The argument uses only scalar unary masks and the coefficient-valued Gram
estimate.  It thus preserves the normalized-trace, fixed-matrix-size, and
homogeneous Hilbert bounds of that estimate.  In the bounded-overlap case,
separate bounded radial and phase masks are already permitted directly.
Finally, ordinary finite-support exhaustion is obtained by unshifted
cutoffs in a proper finite list of logarithmic norms, equal to one on
increasing Cartan balls.  Their supports meet the discrete lattice in
finitely many points.  Their radii may increase as fast as a given test or
approximation requires without changing the bound.
\end{proof}

Choose a vertex tuple $F$ selecting one type in each component, with type
set $I$, and a compact-frame opposite tuple $F^-$.  A tuple path has inverse
split vector
\[
 L=\sum_{i\in I}t_i e_i.
\]
For the height and radial-limit arguments below, take these parameters
successively in any chosen order, starting at one selected vertex and
adding the other selected coordinates singly.  The mixed-face angular
law itself also holds for simultaneous parameters, by
Proposition~\ref{prop:transfer-sharp-flags}.
We also use full-chamber lexicographic paths through $F$, in any order,
with opposite chamber through $F^-$; then $L$ includes all coordinates.
The preceding angular statements and bounded complementary thickness hold
for these paths, by sharpness, exact precomposition, and the total upper
count.  For an actual Fourier index $h$ put
\[
 r_i(h)=\lambda_{p(i)}^T(\mu(h^{-1}))-\lambda_i^S(L).
\]

The next lemma is first applied to the isolated heights of
Proposition~\ref{prop:scale-isolated-height}, yielding the isolated
tightness in Lemma~\ref{lem:heights-tightness-test}. Its later tuple
applications use the heights of Proposition~\ref{prop:heights-tuples}
and the auxiliary tight centered widths supplied by
Proposition~\ref{prop:heights-tuple-calibration}.

\begin{lemma}[Ordered imaginary-power transfer]\label{lem:heights-transfer}
Suppose that $b=b_{F,i}$ is a positive injective operator affiliated with
$\mathcal A_j$, commuting with its own angular measurement, and satisfying
\[
 \sigma_g b_{F,i}=j_S(x_i,g)b_{gF,i},\qquad
 \beta_h b_{F,i}=b_{F,i}j_T(P_{x_i},h^{-1}).
\]
In the isolated case take the height $b_{x_i}$ of
Proposition~\ref{prop:scale-isolated-height} and require only transversality of its end-input
value $v$ at the same positive source label with the reverse output dual
value $w$.  In the tuple case assume that $b$ commutes with $P_F$, require
the corresponding tuple transversality, and assume tight centered widths
on an auxiliary mixed tuple path at $F$.

On buffered transverse patches, for a tuple or full main path and bounded
frequencies $\theta_n\to\theta$ in the stated iterated parameter limit,
\begin{equation}\label{eq:input-15}
 J(1\otimes b^{\mathrm i\theta})
 =E(e^{\mathrm i\theta_n r_i})\,k_T(w,v)^{\mathrm i\theta}
 J(\bar b^{\mathrm i\theta}\otimes1).
\end{equation}
The bar conjugates the positive operator before its imaginary power is
taken.  A joint symbol in this formula means output multiplication followed
by the adjacent end-input multiplication in the indicated order.
\end{lemma}

\begin{proof}
At fixed path parameters, move the start power of frequency $\theta_n$
through $X_a$.  Since $a^{-1}$ has split vector $L$, its source covariance
produces $\exp(-\mathrm i\theta_n\lambda_i^S(L))$.  To read coefficient
$(h,r)$, pair against $\zeta(ay)w_{(h,r)}$, where $r$ comprises the
unmeasured indices.  Moving the end action through this right unitary
changes $\zeta$ to
\[
 b^{-\mathrm i\theta_n}
 j_T(P_{x_i},h^{-1})^{-\mathrm i\theta_n}\zeta.
\]
Conjugation in the first input gives exactly the signs of
\eqref{eq:input-15}.  The remaining group, base, and coefficient multipliers
commute with the power.

On this moved side take a bounded start column before any movement.  The
end column may be Hilbert: the finite-time bilinear estimate with one
bounded input is uniform.  Lemma~\ref{lem:heights-line-normalization} and
uniform separated expansions replace the target factor by the pairing
phase on transverse patches.  Imaginary powers at $\theta_n$ converge in
Hilbert norm on each fixed column to the power at $\theta$.

We justify the original side without asserting a product estimate for two
arbitrary Hilbert columns.  Construct the end vector as a cut one-ended
effect applied to a bounded column $\zeta_0$.  Its defining auxiliary path
is a pure vertex path at $x_i$ in the isolated case, with smooth angular
cut and compact complementary thickness.  In the later tuple case it is a
mixed path at $F$ with bounded centered widths.  Choose its image patch
buffered transverse to the reverse main window.  By sharpness the ranges
of such vectors exhaust the corresponding angular patches.

First retain finitely many measured indices.  Test the main link against
bounded-column arrays, with smooth output tests at fixed thickness and
sufficiently large prescribed gaps; allow powers of the retained phase as
well.  At a fixed main time the cut end vector can be replaced by its
defining auxiliary weak limit.  This remains valid when the other input is
the Hilbert error from replacing a height-applied start vector by a bounded
approximation: at finite measured support the remaining coefficients are
bounded operators, and normality permits bounded tests in the full
remaining tracial algebra.

More concretely, let $f$ be bounded and let $f_l$ be a bounded-column
approximation to $b^{\mathrm i\theta}f$.  After trace rearrangement the
start error is paired against a product of the form
\[
 (D_{\mathrm{aux}}X_s(ay)\zeta_0(ay))
 (D_{\mathrm{main}}T),
\]
where $T$ is the bounded-column testing array.  Its intermediate directions
are $q^{-1}o$ and $ho$.  Transversality to the reverse tuple of a full
chamber suffices for a tuple patch by opposite-face projection.  In the
tuple case the overlap widths are bounded; in the isolated case the
auxiliary support is a single smooth vertex cone, with compact cuts in the
other factors.  Hence \eqref{eq:input-2} and
Lemma~\ref{lem:heights-moving-tests} bound this product uniformly.  The error
is at most that fixed bound times
\[
 \|b^{\mathrm i\theta_n}f-f_l\|_2.
\]
Common projective transport phases disappear from this tracial estimate.
First choose $l$ and then pass to the parameter limit; the error tends to
zero as $l\to\infty$.

The finite-support tests can use the smooth radius cutoffs of
Lemma~\ref{lem:heights-moving-tests}.  A phase may be computed from the
logarithmic Hilbert norm of the weight matrix: at a divergent pertinent gap
the difference from the logarithmic top norm tends to zero by rank one.
At a finite place the top growth is computed exactly on a compact
transverse input line for all sufficiently large gaps, after division by
the bounded locally constant reverse pairing.  Thus all phase/radius tests
retain the uniform estimate before supports are removed.

At each main time the radii may be made sufficiently large for the
moved-side expansions and the bounded start approximant.  Their product
error bound is independent of those radii.  After the bounded start
approximation, the end vector may remain Hilbert by the one-bounded-input
estimate.  Exhaust auxiliary thickness and transverse patches, then the
main stratum using ambient continuous tests with slack and a final Borel
output projection.  Testing in the spectral span generated by these
compact variables and the retained phase proves the vector equality.
Additional phases not used in this equality need not be included in that
span.  A second measured leg can remain throughout in the coefficient
algebra.  This proves \eqref{eq:input-15} with the stated order and domains.
\end{proof}

\begin{lemma}[Total transverse ranges and shrinking frequencies]
\label{lem:heights-tightness-test}
The closed sum of the ranges of transverse ordered angular rectangles for
opposite source tuples is the whole endpoint space.  The same holds for
the isolated vertex-pair rectangles.  In the isolated application of
\eqref{eq:input-15}, the centered variable $r_i$ is tight, including along
the subsequent offset limits.
\end{lemma}

\begin{proof}
The closed range sum is target invariant.  It is also fixed by an
all-factor escaping source subgroup fixing the relevant labels and
normalizing the opposite radical moves.  Its projection is therefore
scalar by the link invariance argument.  It is nonzero: simultaneous source
transport to generic pairs, followed by patch fullness, supplies a nonzero
ordered rectangle.  This proves totality.  The tuple assertion follows by
the same argument on the open opposite-tuple orbit.

For the isolated height put $\theta_n\to0$ in
\eqref{eq:input-15}.  Pull the sharp reverse angular variable through $J$.
The pairing phase and both imaginary powers tend to one on each buffered
rectangle.  Hence the retained phase fixes $J$ on its total range sum and
therefore everywhere.  If tightness failed, some fixed bounded-column input
would retain a positive mass beyond thresholds tending sufficiently slowly
to infinity.  Averaging
$|e^{\mathrm it r}-1|^2$ over a shrinking interval of $t$ gives, by the
elementary sinc formula, a uniformly positive value on that escaped mass.
Choose a frequency in each such interval.  These frequencies tend to zero
and contradict the just-proved equality.  Thresholds and testing vectors
can be selected by countable exhaustion along the successive path
filters.  The same proof permits offsets held after all ordinary times,
so the centered isolated widths are tight on those iterated paths too.
\end{proof}

\subsection{Two-volume calibration of tuple paths}

\begin{lemma}[Coefficient integral bound]\label{lem:heights-coefficient-integral}
Choose a finite Parseval frame $(f_l)$ for the representation of
$\pi_0$ and $W_{R'_j}$ on $H$: the full group translates of these constant
Hilbert columns form a Parseval family.  Let $b_0\ne0$ be a
$W_R$-bounded end column.  Up to the fixed Haar and module normalization,
for a finite set $D\subset\Gamma_j$ the integral over $a\in G$ of the
total measured coefficient probability at $h\in D$, summed over $l$, is
at most
\[
 |D|\,\|b_0\|_2^2.
\]
The path probability comparisons can be integrated over bounded compact
rotation and split-offset parameters by exact pre/post translation of
these inputs, also through successive ultralimits.
\end{lemma}

\begin{proof}
Unfold, for fixed $y$, the identity
$a d_0(y)=d_0(y')\gamma$.  At fixed $h$ the integrated coefficient sum is,
up to the fixed normalization, the sum of
\[
 \big|\langle W_{h,r}^*b_0(y'),\pi_0(\gamma)f_l\rangle\big|^2
\]
over $\gamma,r,l$, at the independent base variables $y,y'$.  Move $W_r$
to the complementary frame index.  Complementary Parseval and commutation
up to scalar phases make the whole sum equal to one squared norm of
$b_0(y')$.  Integration proves the bound.  In particular no extra factor
comes from summing the unmeasured group indices.

A bounded set of compact rotations and permissible split offsets gives a
compact family of translated inputs in Hilbert norm, preserving the
operator bound of the bounded column.  Approximate the Hilbert column
first by bounded ones, and then both translated families by finite compact
nets.  The one-bounded-input probability estimate controls the errors
uniformly in path time.  Exact translations therefore permit passage
under these bounded parameter averages.  Apply the argument successively
for the different path times; no exchange with an arbitrary nonuniform
ultralimit is needed.
\end{proof}

\begin{proposition}[Tuple calibration]\label{prop:heights-tuple-calibration}
For every selected tuple position $i$,
\begin{equation}\label{eq:input-16}
 2\rho_T(s_i e_{p(i)})=2\rho_S(e_i),\qquad
 \mu(h^{-1})=\sum_{i\in I}s_it_i e_{p(i)}+O_{\mathrm{tight}}(1).
\end{equation}
These assertions hold with the total parameters on the further offset
limits specified above.
\end{proposition}

\begin{proof}
We first prove the high-coordinate upper bounds.  Compare a tuple path
$a$ with a faster opposite many-panel lexicographic face path $d$.
All the parameters of $d$ are taken to their limits before those of $a$.
The path $d$ omits exactly
the complementary selected high types $i^*$ and retains full isolated
directions.  Thus $da^{-1}$ has split vector
$L+L_{\mathrm{panel}}$, with the panel part faster.  Cut the two columns to
transverse start angular patches and compact complementary thickness.
Also impose bounded centered isolated widths on the tuple column, using
Lemma~\ref{lem:heights-tightness-test}, and test a high excess
\[
 \alpha_{p(i)}(\mu(h^{-1}))>s_i t_i+M.
\]
Trace cyclicity pairs the cut $X_d$ column times the adjoint of the cut
$X_a$ column against $X_dX_a^*$, which, after base change and a harmless
common phase, is the $da^{-1}$ column.  The product has a uniform Hilbert
bound by the bounded-overlap case of \eqref{eq:input-2}.  Transverse Cartan
addition forces excess greater than $M-C$ in the pertinent residual
quotient gap.  Apply the local upper-speed estimate, \eqref{eq:finite-speed} or
\eqref{eq:input-11}, to the inverse path in the opposite orientation. Its missing-type bound
is exactly the needed bound on the
ordinary index product.  Cauchy--Schwarz therefore makes this pairing tend
to zero as $M\to\infty$.

Exhaust continuity cuts, complementary thickness, and angular restrictions.
If positive infinite excess remained, its one-ended angular sublaw would
be a positive scalar times $P_F$.  Its ordered products with all transverse
opposite many-panel patches could not vanish: simultaneous conjugation of
the source tuple and opposite subfaces ranges over their open orbit, and
fullness supplies nonzero products on conull such choices.  The scalar
mass is unchanged by these fixed conjugations, which only cost bounded
shifts or tight perturbations.  This contradiction proves all high upper
bounds.  The isolated ones already follow from centered isolated
tightness. Both local upper-speed estimates use the total shifted
parameters, so
this reasoning also applies with the permitted further offsets.

Use the frame and bounded end column of
Lemma~\ref{lem:heights-coefficient-integral}.  Around $L$ integrate over a
fixed thick shell, with independent compact rotations and offsets in a
fixed box whose complementary coordinates are positive.  At a real wall
take an interior subbox; the Cartan density then gives the lower bound
\[
 c\exp(2\rho_S(L)).
\]
At a finite place integrate over the disjoint selected Cartan double cosets
in the permissible mesh.  Each is bi-compact averaging times a constant
comparable to $\exp(2\rho_S(L+\eta))$, by the root-coordinate index count.
The same lower estimate follows.  The mixed link isometricity, the just
proved upper bounds, and complementary thickness retain a positive mass
after sufficiently large fixed cuts.  This can be done uniformly over the
bounded translation parameters by the compact-net argument of the lemma.

Disjoint fixed neighborhood thickening of lattice points and Cartan
integration bound the target count in the corresponding windows by
\[
 C\exp\!\left(\sum_{i\in I}t_i\,2\rho_T(s_i e_{p(i)})\right).
\]
At finite places include all Cartan entries in the window enlarged by a
fixed constant; their multiplicity per unit mesh box is bounded.  The
coefficient integral bound gives the comparison of these exponentials.
Make $i$ the fastest coordinate to obtain
\[
 2\rho_S(e_i)\leq2\rho_T(s_i e_{p(i)}).
\]
The lists are two copies of the same type list, and the speed ratios
telescope around each permutation cycle.  Multiplication around that
cycle gives equality at every position.

Suppose next that a selected coordinate has positive mass of infinite
deficit from its calibrated center. This sublaw again has scalar mass on
the mixed link. Fix complementary and other upper widths large enough to
retain a mass $\varepsilon>0$, and then impose deficit greater than $M$.
The same lower shell estimate and the coefficient integral bound give
\[
 \varepsilon c\exp(2\rho_S(L))
 \le C\exp(2\rho_S(L))
         \exp\bigl(-M\,2\rho_T(e_{p(i)})\bigr).
\]
Here the cut widths are fixed before $M$ is exhausted, so $c,C>0$ do
not depend on $M$. The final factor is the gain in the target packing
bound. In simple coordinates every coefficient $2\rho_T(e_j)$ is
positive, and integration up to each upper bound costs a constant times
its endpoint exponential. There is consequently no polynomial length
loss. The displayed inequality is impossible for large $M$, proving
tightness of every deficit.

The same argument applies to any purported positive mass in a subsequent
offset limit: offsets are fixed during all ordinary path times and are
exhausted only afterwards. Finally, $p$ preserves the Coxeter components
by Section~\ref{sec:transfer}, and the tuple selects only one type in each
component. Evaluating $\lambda_{p(i)}^T$ in
\eqref{eq:input-16}, using
$\lambda_{p(i)}^T(e_{p(i)})=2\rho_T(e_{p(i)})$, therefore proves tightness
of each $r_i$. This evaluation requires no cross-weight Cartan identity.
\end{proof}

\subsection{Height solution lines and pair identities}

\begin{proposition}[Tuple heights]\label{prop:heights-tuples}
At each non-isolated tuple position there is a positive injective affiliated
height $b_{F,i}$ commuting with $P_F$ and satisfying
\begin{equation}\label{eq:input-17}
 \sigma_g b_{F,i}=j_S(x_i,g)b_{gF,i},\qquad
 \beta_h b_{F,i}=b_{F,i}j_T(P_{x_i},h^{-1}).
\end{equation}
At an isolated position the previous height $b_{x_i}$ can be used and also
commutes with $P_F$.  On the tuple link,
\begin{equation}\label{eq:input-18}
 J^*E(e^{\mathrm i\theta r_i})J
 =\chi(\theta)
   \bar b_{F^-,i^*}^{\mathrm i\theta}\otimes b_{F,i}^{\mathrm i\theta},
\end{equation}
where $\chi$ is a scalar probability characteristic function.  Heights may
be chosen in the two-slot system and then copied normally to the system
with two measured legs.
\end{proposition}

\begin{proof}
The compressed characteristic operators of the tight variable $r_i$ commute
with the sharp tuple coordinates.  They lie in
$\overline{\mathcal A_j}\bar\otimes\mathcal A_j$, have independent
contracted invariances, and transform under independent target shifts by
the endpoint line phases of Lemma~\ref{lem:heights-line-normalization}.
At a given end and frequency $t$, a bounded solution $B_t$, commuting with
the pertinent sharp vertex, has the tilted law
\[
 \beta_h B_t=B_tj_T(P_{x_i},h^{-1})^{\mathrm it}.
\]
For two solutions of the same frequency, both orders of their products
with adjoints have untilted target laws and the contracted invariances.
They are in $\mathcal N$ and scalar.  A nonzero solution is consequently
a scalar multiple of a unitary, and the solution space has dimension at
most one.

Normal slices of the characteristic operators give locally nonzero
solutions near zero, strongly continuous together with their adjoints.
Here tightness supplies strong convergence of the characteristic operators
at zero.  Normalize the nonzero slices to unitaries.  Products of these
local solutions give the corresponding solution line at every frequency,
and the lines form a continuous projective one-parameter unitary group.
They commute with $P_F$ because the original compressed operators do.

The group consisting of those unitary solution lines has local continuous
sections over $\mathbb R$ and circle fibers, hence is locally a
two-dimensional Euclidean group.  The locally Euclidean groups theorem
makes it a Lie group \cite{Gleason1952,MontgomeryZippin1952}. Its Lie-algebra extension over the one-dimensional
abelian Lie algebra splits; integrating that section gives a continuous
splitting over the simply connected real line.  Stone's theorem \cite[Theorem~B]{Stone1932} now writes
the resulting group as $b_{F,i}^{\mathrm it}$ for the exponential of a
selfadjoint affiliated operator.  Thus $b_{F,i}$ is positive and injective.
One-dimensionality of both end solution lines gives
\eqref{eq:input-18}.  After the genuine unitary groups are removed, its
scalar coefficient is continuous and positive definite: in the
operator-valued positive-definite matrix tests absorb the inverse unitary
group factors into the testing vectors.  Its value at zero is one.

We determine the source covariance rather than assuming its character.
The source parabolic stabilizer of $F$ normalizes these intrinsic solution
laws and therefore acts on the height by a continuous positive character.
Replace $L$ by $L+mA_0$ along a selected split-center direction, with the
integer $m$ held after all ordinary path times.  Before recentering,
exact precomposition translates the characteristic law by
$m\log(\operatorname{char}(A_0))$ in the start slot.  Recentering subtracts
$m\lambda_i^S(A_0)$.  Proposition~\ref{prop:heights-tuple-calibration}
gives tightness even when these subsequent multiples are exhausted.
A nonzero residual drift would contradict that tightness on each unit
input.  Hence the two characters agree on every selected split direction.
Those directions span the split center of the tuple parabolic.  A real
logarithmic character vanishes on its radical, compact parts, and derived
Levi factors, as explained above; the finite-place mesh determines it
there as well.  Compact transport now gives exactly
\eqref{eq:input-17} at all labels.

At an isolated position retain the previously constructed $b_{x_i}$.
It solves the same one-dimensional laws, so its powers agree up to scalars
with the tuple solution and commute with $P_F$.  The reverse-end heights
solve the same intrinsic laws at opposite positions.  All constructions
use countable determining labels and tests before taking Borel output
cuts; strong continuity then supplies every fixed label.  The normal-copy
argument of Section~\ref{sec:links}, together with uniqueness of the
solution lines, therefore permits these one-ended heights to be copied
to the two-measured-leg system.  Exact end-input identities are copied
normally too.  Later commutation comparisons use one common all-gap path
in that system, rather than unrelated marginal limits.
\end{proof}

\begin{proposition}[Ordered pair law]\label{prop:heights-pair-law}
For opposite tuples $Y,F$, on transverse tuple patches,
\begin{equation}\label{eq:input-19}
 \begin{split}
 b_{Y,i^*}^{\mathrm i\theta}P_Y(dw)\,
 b_{F,i}^{\mathrm i\theta}P_F(dv)\,
 k_T(w,v)^{\mathrm i\theta}
 &=C(Y,F)^{\mathrm i\theta}P_Y(dw)P_F(dv),\\
 C(Y,F)&=c\,k_S(y_{i^*},x_i),\qquad c>0.
 \end{split}
\end{equation}
The statement also holds on the conjugate slot, with conjugation and the
corresponding sign change.  At an isolated position only the vertex-pair
transversality is needed.
\end{proposition}

\begin{proof}
Apply \eqref{eq:input-15} to a tuple path at constant frequency, and compress
using \eqref{eq:input-18}.  Cancel the independent start power and pull the
sharp reverse tuple through $J$.  At $Y=F^-$ this gives the left side of
\eqref{eq:input-19}, multiplied by $\chi(\theta)$, equal to the unweighted
ordered product.

This equality can be compared at both signs despite the ordering.  Each
height power commutes with its own spectral slot.  Multiply outside the
left and right slots by their inverse powers and insert the nonvanishing
smooth pairing phase on the transverse patch.  Applying the opposite-sign
identity gives
\[
 \chi(\theta)\chi(-\theta)\,P_Y(dw)P_F(dv)
 =P_Y(dw)P_F(dv)
\]
on such patches.  Nonzero ordered rectangles exist by
Lemma~\ref{lem:heights-tightness-test}; hence $|\chi(\theta)|=1$.
Since $\chi$ is a probability characteristic function, it is the
characteristic function of a point mass.  Its scalar is a positive
constant raised to the imaginary power.  Simultaneous source transport and
the pairing covariance give $C(Y,F)=c k_S(y_{i^*},x_i)$.

For clarity about subsequent substitutions, write $P,Q$ for the reverse
and same-label tuple slots on a conjugate end input, and $A_t,B_t$ for the
corresponding conjugated height powers.  The compressed equation reads
\[
 \chi(t)A_t P B_t Q\,k_T^{\mathrm it}=PQ.
\]
Thus the substitution used on a full path is
$P B_tQ\,k_T^{\mathrm it}=C^{\mathrm it}A_t^{-1}PQ$, in exactly that
order.  The relevant total range sum may be multiplied on the left by
the inverse reverse power and remains total.  Smooth multipliers and
separated expansions on the buffered patches justify every localized
identity.  The isolated case follows from the first version of
\eqref{eq:input-15} in the same way.
\end{proof}

\subsection{The full centered law}

\begin{proposition}[Sharp full-path heights]\label{prop:heights-full-law}
On a full lexicographic path,
\begin{equation}\label{eq:input-20}
 E(e^{\mathrm i\theta_n r_i})J
 =J\!\left(C(F^-,F)^{-\mathrm i\theta}
   \bar b_{F^-,i^*}^{\mathrm i\theta}\otimes b_{F,i}^{\mathrm i\theta}\right)
 \qquad(\theta_n\to\theta).
\end{equation}
There is an invertible linear map $A_*$ on the split-coordinate spaces
satisfying
\begin{equation}\label{eq:heights-weight-map}
 \lambda_{p(i)}^T A_*=\lambda_i^S,\qquad
 A_*e_i=s_i e_{p(i)},\qquad 2\rho_TA_*=2\rho_S.
\end{equation}
The centered Cartan vector $\Delta=\mu(h^{-1})-A_*L$ has a tight sharp
commuting law on input.
\end{proposition}

\begin{proof}
Insert \eqref{eq:input-19} into \eqref{eq:input-15} on a full path.
Pulling back the sharp reverse tuple first produces the ordered
conjugate-end expression
$\bar P_{F^-}(dw)\bar b_{F,i}^{\mathrm i\theta}\bar P_F(dv)$.
The ordered substitution in the proof of
Proposition~\ref{prop:heights-pair-law} gives
\eqref{eq:input-20} on transverse ordered end ranges acted on by the inverse
reverse power.  Those ranges are total.  Letting frequencies tend to zero
and applying the shrinking-frequency argument proves tightness on the
full path too.

Apply this for every type using any tuples in the chamber.  The determinant
weights form bases, so the first equation in
\eqref{eq:heights-weight-map} determines an invertible $A_*$.  To identify
its column at $e_i$, choose a tuple prefix with $i$ fastest, add its
other selected types at successively slower scales, and then add the
remaining chamber types one at a time.  Dividing the centered law by that fastest time and
using \eqref{eq:input-16}, with the subsequent translations handled by exact
precomposition, gives $A_*e_i=s_i e_{p(i)}$.  The last equality follows
from the calibrated $2\rho$ values in \eqref{eq:input-16}.  The output
Cartan phases commute; \eqref{eq:input-20} pulls their sharp law through the
isometry, giving the asserted joint law of $\Delta$.
\end{proof}

The full-path law also controls the height groups before arbitrary-label
commutation has been proved. Every height belonging to a chamber commutes
with that chamber's angular measurements. This includes the angular
measurements in either normal measured copy: they and the centered output
phases belong to the same full-path spectral law, and
\eqref{eq:input-20} pulls them through the same isometry. More explicitly,
write $\bar U_-(t)\otimes U_+(t)$ for the two endpoint height groups
in that equation, with its scalar factor suppressed. If $P$ is a chamber
angular projection in either measured copy at the start input, then
\[
 [\bar U_-(t)\otimes U_+(t),1\otimes P]
   =\bar U_-(t)\otimes[U_+(t),P]=0.
\]
The reverse factor is unitary, so $[U_+(t),P]=0$.

Two height groups belonging to that chamber may at this stage commute
only up to a scalar. Indeed their tensor products in
\eqref{eq:input-20} commute, so tensor cancellation makes the component
commutators reciprocal scalars. The one-parameter group laws make each
scalar commutator a bicharacter,
and strong continuity gives the form $e^{\mathrm icst}$ in the two real
frequencies.
The same argument applies across measured copies because both were
retained in the common full-path law.

For the two copies of the same isolated height $b_x$, this scalar is one.
On $H=K\otimes K\otimes K$ from Section~\ref{sec:setup}, use the
complex-linear unitary
\[
 \mathcal F(\xi_0\otimes\xi_1\otimes\xi_2)
   =\xi_0\otimes\xi_2\otimes\xi_1.
\]
It fixes the trunk action and diagonal $\rho$, interchanges the two
measured actions, and preserves the multiplier
\eqref{eq:setup-tensor-multiplier}. It acts fibrewise on the induced
space. Normal copying in Proposition~\ref{prop:links-link} and the
intrinsic closed derivative modulus in
Proposition~\ref{prop:scale-isolated-height} therefore give
$\mathcal F b_x^1\mathcal F^*=b_x^2$ and the converse equality.
If $U_j(t)=(b_x^j)^{\mathrm it}$ and
\[
 U_1(s)U_2(t)=e^{\mathrm icst}U_2(t)U_1(s),
\]
conjugation by $\mathcal F$ and comparison with the relation with $s,t$
interchanged give $e^{2\mathrm icst}=1$ for all $s,t\in\mathbb R$.
Thus $c=0$. The two copies of the complete same-label angle-height
calculus consequently commute strongly. This interchange of measured
slots is separate from the conjugate endpoint factor in the full-path
law.

Write
\[
 2\rho_S=\sum_i a_i\lambda_i^S.
\]
Equation~\eqref{eq:heights-weight-map} gives the same coefficients for the
permuted target weights.  At a chamber $Z$ fix, by their type sets, tuple
choices $F(i)\subset Z$ containing $i$.  Define
\[
 D_Z^{\mathrm it}=
 \left(\prod_i b_{F(i),i}^{\mathrm it a_i/2}\right)
 \left(\prod_i^{\mathrm{reverse}} b_{F(i),i}^{\mathrm it a_i/2}\right).
\]
The symmetric order cancels all scalar Weyl reordering phases, so this is
a genuine strongly continuous unitary group.  It defines a positive
injective affiliated height $D_Z$.  Use opposite tuple choices to define
$D'_{Z'}$ at the other end.  These heights commute with their chamber
angles and have the covariance laws \eqref{eq:input-17} with chamber
cocycles.  On independently compact-rotated full paths, spectral calculus
in \eqref{eq:input-20} gives
\begin{equation}\label{eq:input-21}
 \exp(2\rho_T(\Delta))
 =\operatorname{const}\cdot\bar D'_{Z'}\otimes D_Z,
 \qquad \operatorname{const}>0.
\end{equation}
Compact rotations are exact input translations and preserve the line
norms separately.  A further fixed shift $L\mapsto L+\eta$ shifts the
law relative to the old center by $A_*\eta$.  These statements concern
one common joint calculus, not separately chosen height marginals.

\subsection{Weighted chamber probability}

\begin{theorem}[Haar-height isometry]\label{thm:heights-isometry}
Let $B_S,B_T$ be the source and target chamber spaces with their compact
probabilities.  There is a constant $0<c_0<\infty$ and an isometry
\begin{equation}\label{eq:input-22}
 \begin{split}
 V:L^2(B_T;\mathcal H)&\longrightarrow L^2(B_S;\mathcal H),\\
 (V(f\otimes\xi))(Z)&=c_0^{-1/2}D_Z^{-1/2}P_Z(f)\xi.
 \end{split}
\end{equation}
The expression on the second line is initially interpreted through positive
form integrals and then extended by the stated isometry.
\end{theorem}

\begin{proof}
\emph{The weighted packing estimate.}
Count inverse lattice indices in a fixed-width box around $A_*L$ whose
image chamber lies in a continuity patch $O\subset B_T$.  Pack with fixed
disjoint right neighborhoods of those inverse indices.  In all-gap
windows, right thickening perturbs the image chamber by $o(1)$ uniformly,
by highest-line convergence, and changes the Cartan vector by a bounded
amount.  Cartan integration therefore gives
\[
 \limsup\frac{\#\{\text{such indices}\}}{\exp(2\rho_S(L))}
 \leq C|O|.
\]
One may first enlarge angular patches by fixed slack and then shrink it by
compact regular approximation.  The finite-place count includes all
double cosets in a fixed enlargement of the box, with bounded multiplicity
per mesh box, as in Proposition~\ref{prop:heights-tuple-calibration}.

Integrate both compact rotations and split shifts $\eta$ in a bounded box,
using the frame inputs $\bar b_0\otimes f_l$ of
Lemma~\ref{lem:heights-coefficient-integral}.  The coefficient bound and
the Cartan density, normalized by $\exp(2\rho_S(L))$, show that a constant
times $|O|$ bounds the weighted integral of limiting probabilities of
\[
 \Delta+A_*\eta\ \text{in the fixed box},
 \qquad\text{image chamber in }O,
\]
with weight $\exp(2\rho_S(\eta))$.  Begin with smaller interiors and
continuous tests.  Exact input translations and the compact-net argument
permit the bounded parameter integral before each limit.  Equation
\eqref{eq:input-20} identifies these translated tests with the same sharp
law, shifted by $A_*\eta$ about the old center; compact rotations introduce
no extra line scalings.

Now exhaust the whole split space, using a sum on finite-place meshes and
ordinary integration on real coordinates.  Choose the fixed counting box
sufficiently wide.  For every $\Delta$, the set of permissible offsets for
which $\Delta+A_*\eta$ lies in a smaller box has uniformly positive
sum-volume near $-A_*^{-1}\Delta$.  Since $A_*$ is invertible and
$2\rho_TA_*=2\rho_S$, this yields
\[
 \int_{\{\eta:\,\Delta+A_*\eta\ \text{in the box}\}}
       e^{2\rho_S(\eta)}\,d\eta
 \geq c e^{-2\rho_T(\Delta)}.
\]
Here and below this notation includes the indicated mesh sums.  A full
lattice has a bounded covering radius, which proves the assertion at
finite places; the real subbox has a fixed positive volume.  Only bounded
offset windows are used before a path limit, so all chambers remain deep
in the prescribed iterated order.

Positive integration thus bounds the expectation of
$e^{-2\rho_T(\Delta)}$ with the angular cut.  By
\eqref{eq:input-21}, independence of the two compact rotations, and the
product input $\bar b_0\otimes f_l$, that positive form factors into an
opposite factor and a forward factor, up to a fixed constant.  No
exponential moments have been presumed.  Perform this step with positive
spectral truncations first; any strictly positive bounded truncation of
the averaged opposite factor is enough to divide.  It follows that
\begin{equation}\label{eq:heights-weighted-bound}
 \sum_l\int_{B_S}
 \langle f_l,D_Z^{-1}P_Z(O)f_l\rangle\,dZ
 \leq C'|O|.
\end{equation}
Regular approximation extends the bound from continuity patches to all
positive measurable tests.

\emph{The positive density map.}
For bounded $f\geq0$ define the extended positive form
\[
 T(f)=\int_{B_S}D_Z^{-1}P_Z(f)\,dZ.
\]
The height commutes with its angular slot, so this integral is positive.
Its source covariance and the chamber Jacobian place it in the extended
positive part of $(\mathcal A_j)^\sigma$.  Indeed the inverse line factor
is canceled by the change of angular probability under the source
transport.  After undoing induction a $\sigma$-fixed field is constant
with value in the complementary commutant.  Evaluation on the finite
Parseval frame, summed over $l$, is precisely its faithful finite trace.
Thus \eqref{eq:heights-weighted-bound} makes $T$ a bounded positive map
\[
 L^1(B_T)\longrightarrow L^1((\mathcal A_j)^\sigma).
\]
It intertwines transport of chamber densities with the trace-preserving
target action.  This also justifies the claimed $L^1$ values of the
unbounded positive form integral.

\emph{Why this map is scalar.}
Induce the finite algebra $(\mathcal A_j)^\sigma$ along the target lattice
inclusion, with its induced finite trace.  If $\mathcal D_g$ denotes
transport of chamber densities, the induced equivariant map is represented
by fields
\[
 g\longmapsto T(\mathcal D_{g^{-1}}f).
\]
Their equivariance under the right lattice action follows from that of
$T$.  They are integrable over the quotient: $\mathcal D_{g^{-1}}$ is a
positive $L^1$ isometry, and the norm of $T$ gives a uniform bound on their
$L^1$ norms.  Let $h_0$ be the induced image of the constant density $1$,
and let $\alpha$ be the induced group action.

For every target group element $s$, averaging $\mathcal D_k\mathcal D_s1$
over the special compact group gives $1$, because that compact group is
transitive on chambers and the density has total mass one.  Equivariance
therefore gives
\[
 h_0=\int_K\alpha_k\alpha_s(h_0)\,dk.
\]
Every integrand has the same spectral distribution as $h_0$, since the
induced action preserves the finite trace.  For a threshold $t$ which is
not an atom of that distribution, the projection
$p_t=1_{(t,\infty)}(h_0)$ uniquely maximizes
$\tau(p(h_0-t))$ among projections.  The maximum is the same for each
integrand.  Testing the displayed average against $p_t$ gives equality in
the average of these upper bounds, so $p_t$ is the maximizing spectral
projection for almost every integrand.  Take a countable dense set of
atom-free thresholds; their spectral projections determine the whole
positive $L^1$ element.  Hence
$\alpha_k\alpha_s(h_0)=h_0$ for almost every $k$.  Strong $L^1$ continuity
extends this to all $k$, and in particular $\alpha_s(h_0)=h_0$.

The span of the group transports of $1$ is dense in $L^1(B_T)$.  To see
this directly, take a regular split contraction.  It pushes chamber
probability to a point mass away from its null repelling set.  Convolution
of these concentrating densities with continuous compact-rotation weights
approximates continuous chamber densities uniformly, by lifting to the
compact group modulo the chamber stabilizer.  Continuous densities are
dense in $L^1$.  The induced map therefore depends only on total integral.
Take a common conull section value for a countable dense set of density
tests.  Because density transport is an onto $L^1$ isometry, this
rank-one conclusion descends to the original map.  Its value is fixed by
both source and target actions, hence scalar.  We obtain
\[
 T(f)=c_0\int_{B_T} f(v)\,dv.
\]
The trace bound makes $c_0$ finite, and positivity together with the
injectivity of the heights makes it strictly positive.

Finally, for bounded angular tests, commutation of $D_Z$ with $P_Z$ gives
the norm identity in \eqref{eq:input-22} from this formula for $T$.
Polarization handles finite sums of elementary tensors, and completion
gives the claimed isometry.  Positive truncation and separability give the
measurable $L^2$ representative in that equation, without imposing an
operator bound on $D_Z^{-1/2}$.
\end{proof}

\subsection{The height intertwiner}

\begin{proposition}[Intertwining with ordinary chamber variables]
\label{prop:heights-intertwiner}
Fix $F,i$.  For generic source chamber $Z$, let $Y(Z)\subset Z$ have the
tuple type opposite to $F$.  If $v$ is an ordinary target chamber, with
the relevant dual vertex used in the pairing, then
\begin{equation}\label{eq:input-23}
 V\bigl[k_T(v,P_{x_i})^{\mathrm i\theta}
              b_{F,i}^{\mathrm i\theta}\bigr]
 =\bigl[C(Y(Z),F)^{\mathrm i\theta}
              b_{Y(Z),i^*}^{-\mathrm i\theta}\bigr]V.
\end{equation}
Both bracketed operators are unitary.  Moreover $D_Z$ strongly commutes
with each such own-coordinate height $b_{Y(Z),i^*}$.
\end{proposition}

\begin{proof}
For any fixed spectral input value the nontransverse ordinary chamber set
is null.  Fubini therefore gives transversality on the full product of
ordinary chamber probability and the input spectral measure.  The powers
in the two brackets have modulus one and commute with their own angular
slots, proving unitarity.

Localize in the ordinary chamber variable and in the input $P_F$
coordinates on buffered transverse rectangles; these give total domains.
Before the factor $D_Z^{-1/2}$ in $V$, substitute at $P_Z$ using
\eqref{eq:input-19}.  Chamber tests may be placed on the left because the
$Y(Z)$ height commutes with them.  Use separated smooth expansions, or
finite clopen partitions at finite places, for the pairing tests, and
truncate the $D_Z$ height to bounded intervals.

The scalar Weyl relations obtained from \eqref{eq:input-20} show that
commuting $b_{Y(Z),i^*}^{-\mathrm i\theta}$ through $D_Z^{-1/2}$ costs at
most a positive scalar, while its height truncations are only dilated.
This scalar is constant on the generic chamber orbit: simultaneous source
transport preserves the scalar commutator, and the tuple choices inside
a chamber have fixed type sets.  Remove the truncations using the
isometry \eqref{eq:input-22}.  We obtain \eqref{eq:input-23} initially up
to that positive scalar.  Both bracketed operators are unitary and $V$
is an isometry, so comparison of norms forces the scalar to equal one.
For every real frequency this makes the Weyl coefficient vanish; hence
the corresponding unitary groups, and therefore the positive affiliated
heights, commute strongly.  The same argument proves
\eqref{eq:input-23} exactly on the completed domains.
\end{proof}
